\documentclass{article}
\usepackage{microtype}
\usepackage{graphicx}
\usepackage{subfigure}
\usepackage{booktabs} % for professional tables
\usepackage{multirow}
\usepackage{amsmath}
\usepackage{hyperref}
\newcommand{\theHalgorithm}{\arabic{algorithm}}

\usepackage[accepted]{icml2021}

\icmltitlerunning{Westward Trail: Scrutinizing a Local LLM Game Master}
\begin{document}

\twocolumn[
\icmltitle{Westward Trail: Developing and Scrutinizing a\\Local LLM Game Master via Prompt Engineering}

\begin{icmlauthorlist}
\icmlauthor{Sander De Bleecker (Student number 852656007)}{ou}
\end{icmlauthorlist}

\icmlaffiliation{ou}{Open Universiteit}
\icmlcorrespondingauthor{}{}

\icmlkeywords{prompt engineering, large language models, evaluation, interactive fiction}
\vskip 0.3in
]

\begin{abstract}
I developed an Oregon-Trail-style survival game \emph{Westward Trail}, in which a locally
served open-weight LLM acts as game master: it narrates events and attaches numeric
consequences to the player's choices.
The program itself still holds the authoritative game state and validates every
LLM-proposed effect against hard rules, so the game can never enter an invalid state
\emph{and} every turn doubles as a labeled test case.
On top of this app I run controlled scrutiny experiments comparing three strategies
along four metrics: JSON format reliability, hard-rule adherence, state tracking over
long conversations (with and without ground-truth state in the prompt), and recall of
facts stated once at game start.
Format compliance turns out to be nearly solved even by a minimal prompt (99.5\% parse
rate), but rule adherence is not: a single worked example cuts hard-rule violations from
185 to 14 per 100 turns, while writing the rules out verbatim helps not at all, because
the dominant error is a schema misreading (absolute values echoed into a delta field)
that examples fix and rules cannot.
Without ground-truth state in the prompt, the model's believed food count drifts by over
100 lbs within 20 turns, and recall of facts stated once collapses from 100\% to 6\% at
the moment the game intro slides out of the rolling context window; the model was never
remembering, it was reading.
\end{abstract}

\section{Introduction}

The first idea was to create a roguelike survival game heavily inspired by Oregon Trail,
in which a family tries to survive a multi-month journey with a cart through an
unforgiving wasteland.
There is a dark side to the game where it often just goes wrong for the company and you
have to replay it multiple times until luck is on your side.

This concept was appealing to me, as the chances must have been predefined to have
certain events occur during the playthrough.
Now, with how easy it is to run LLMs locally, the chances can be made nondeterministic or
unknown by simply using the weights of an LLM.
At first all chances were done by the LLM itself; later on the game had to be playable, so
some chances were left to the program.
The two mechanisms, prompted event results versus purely stochastic event
results, are cleanly separated by giving the pure stochastic chances their own mechanic,
namely ``scavenging'': a mechanic where the player can choose to scavenge for items at
night, which has a chance of finding items but also a chance of getting hurt or losing
items.

% Prompted LLMs can be turned into single-purpose applications without any training, but the reliability of such applications is an empirical question: the model may ignore the output schema, break stated rules, or silently lose track of context. We study this in a task that stresses all three at once: an LLM \emph{game master} for a clone of the 1971 educational game \emph{The Oregon Trail}. Each turn the model must (i) produce schema-valid JSON, (ii) invent a short narrative and three options, (iii) attach numeric resource deltas that respect hard game rules, and (iv) stay consistent with everything that happened earlier in a growing conversation.

% The central design idea is that the application and its evaluation are the same artifact. A Python harness owns the true game state and applies every effect itself; the LLM only ever \emph{proposes}. Because ground truth is always available in the harness, discrepancies between what the model believes, proposes, or remembers and what is actually true are measurable on every single turn, with no external dataset or human annotation. Concretely, we ask:

% \begin{enumerate}
%     \item \textbf{Format reliability:} Does the model reliably produce schema-valid JSON, and does a worked example or explicit rules change that?
%     \item \textbf{Rule adherence:} How often, and in which ways, do proposed effects violate hard game rules?
%     \item \textbf{State tracking:} Can the model reconstruct the numeric game state from the dialog history alone (\emph{blind}), compared to having the true state re-injected every turn (\emph{guided})?
%     \item \textbf{Memory:} Does the model recall party facts stated once at game start when quizzed 5, 10, 15, and 20 turns later?
% \end{enumerate}

% All inference is local (PyTorch + HuggingFace \texttt{transformers}), so the test scripts can hammer the model without usage restrictions, and the whole pipeline is reproducible from the submitted code.

\section{Methods}

\subsection{Game loop}

The architecture was laid out around how the few game actions you can take were supposed
to be handled.

\emph{A day and its decision:} every turn in the game is a day, which starts with a
description of how the day is going and with 3 options for how it should be further
planned out. The game then waits for the user's decision.
\emph{The end of day and its consequences:} with the context of the chosen option, the
LLM makes up a JSON object of rewards or penalties the travelling party gets.
\emph{The optional scavenging:} on the left, during the user's decision, the user can
choose to scavenge for items at night. This is a stochastic mechanic with real stakes.

Here, the whole day including consequences used to be generated by one prompt; this
changed as the accuracy was not up to par.
Later on, the day was split into multiple prompts (two for the narration, to make it more
story-like by switching prompt), while the consequences were split into a separate prompt
giving complete context so the prompt could focus on the consequences.
Furthermore, a check is done to ensure the consequences are valid.

You can find the prompts in the prompts folder of the GitHub repository; their usage can
be found in \texttt{prompts.py}.
An HTML website is loaded as the interface for the user via \texttt{serve.py} and
\texttt{index.html}.

Asynchronously, while the user is still deciding, the LLM is already generating a scene
image for the day; this is done in the background and does not block the user from making
a decision.

\subsection{Inference engine and model choice}

For the LLM, everything was run locally on Ollama at first.
Due to project requirements, the same model used in Ollama was then run via PyTorch.

The preferred model for the game loop is the instruction-tuned Gemma E4B checkpoint
\citep{gemmateam2024gemma}.
Beforehand, different variations of Qwen and Mistral were tested.
Gemma E4B was chosen as it seemed more manipulable via prompt, and it was more efficient
on the NVIDIA GeForce RTX 5080 due to being NF4-quantized.
Trial and error found that Qwen2 often needed more retries for outputting valid JSON, that
Qwen3 models were an improvement but heavier, and that Mistral had more hallucinations.

\subsection{Gradual improvement via prompt engineering}
\label{sec:improvements}

The app did not arrive at its current shape in one pass: the development journal
(\texttt{journey.md}) records sixteen small stages, and nearly every one of them is a
prompt intervention rather than an algorithmic change. They group into four recurring
lessons.

\emph{Pick a model the prompt can steer.} The very first obstacle was the model, not the
prompt: a reasoning-tuned Qwen3 checkpoint intermingled its chain-of-thought with the
expected JSON, and disabling reasoning did not fully stop the leakage, which led to a
curated catalog of non-thinking models as the single source of truth for model selection.
The default context limit also proved too short to keep games interesting beyond ten
turns, so the context window was raised explicitly and treated as a budget from then on
(the rolling history window of Section~\ref{sec:strategies} follows from this).

\emph{Give every prompt one job.} The largest quality jumps came from decomposition. When
one call produced both the options and their rewards, the numbers often contradicted the
text (saving an ox by resting would \emph{cost} an ox), which motivated the
narrator/scorer split of Section~\ref{sec:split}. The same principle was applied twice
more: encounter days with strangers got their own prompt, separate from ordinary travel
days, so the two event kinds stop blending into a monotone middle ground; and a dedicated
post-decision prompt narrates how the chosen option played out; initially the
consequences were visible on the choice itself, which was convenient for debugging but
ruined the immersion of finding out how the day went.

\emph{Prompts are code and rot like code.} Midway through, all prompts were moved out of
Python strings into plain text files, which made them readable, diffable, and cheap to
iterate on; one such iteration uncovered a hand-made bug where the worked JSON example
inside a prompt was itself invalid JSON: the model had been imitating a broken template.
A later pass rewrote the surface tone of every prompt to be consistently western-flavored,
an improvement that would have been tedious while prompts lived inline. The subtlest bug
was context leakage: the fixed 2-bullet cost of night scavenging, present in the prompt's
state summary, kept surfacing in generated stories as bullets the party ``found''; the fix
was to remove state the prompt did not need, i.e.\ prompt \emph{minimization} rather than
addition.

\emph{Know what not to prompt.} Playtesting showed the game could be won by simply forcing
maximum effort every day, and no phrasing made the model punish this reliably. The risk of
serious disease when pushing too hard was therefore moved into the harness as a hard-coded
stochastic rule the model only narrates, the origin of the prompted-versus-stochastic
separation described in the introduction. Where the model repeatedly failed at a job, the
improvement was to take the job away from it.

There are of course more, smaller but noteworthy actions, such as setting a tone for the
output by defining that tone in the prompt, giving an example of the output, and stripping
whole concepts from a prompt when they sidetrack it from its main objective.

\subsection{Prompt strategies}
\label{sec:strategies}
We compare three strategies that each vary exactly one dimension over the previous, so
effects attribute cleanly (Table~\ref{tab:strategies}). All three request the same JSON
object: a \texttt{believed\_state} (the model's own reconstruction of the current state,
used for the state-tracking experiment), a 2--3 sentence \texttt{narrative}, and exactly
three \texttt{options}, each with an \texttt{effects} dict of numeric deltas.

\begin{table}[t]
\caption{Prompt strategies. Each adds one ingredient.}
\label{tab:strategies}
\centering
\small
\begin{tabular}{lp{4.6cm}}
\toprule
Strategy & Adds over the previous \\
\midrule
\texttt{minimal} & Task description + bare JSON schema \\
\texttt{fewshot\_schema} & + one fully worked example turn \\
\texttt{rules\_explicit} & + the hard game rules written out verbatim \\
\bottomrule
\end{tabular}
\end{table}

\subsection{Splitting narration from scoring}
\label{sec:split}
At first the app revealed contradictions between the text of the option and the
consequence paired to it.
For example, a ``brief rest'' costing an ox, or $+105$ health for an unhurt character.
The program therefore splits each turn into a \emph{narrator} call (temperature 0.8) and a
\emph{scorer} call (temperature 0.1, given the true state, the narrative, and the three
option texts, and asked only for effects under the hard rules plus explicit coherence
constraints).
A third call narrates the outcome after the game has applied the chosen effects.
A \texttt{--single-call} flag restores the one-call behavior for A/B comparison.
The scripted experiments deliberately use the one-call strategies of
Table~\ref{tab:strategies}, as they scrutinize the program's behavior.

\section{Experimental Setup}
There is no external dataset: the game itself is the test set, and the harness state is the
ground truth (a deliberate consequence of the assignment's local, reproducible setup). A
scripted bot plays full games of up to 20 turns for every combination of strategy $\times$
mode (\emph{guided}/\emph{blind}) $\times$ seed, choosing uniformly among the offered
options. Every turn is appended as one JSON line with a deduplication key, so runs are
resumable and a Ctrl-C never corrupts data; completed runs are skipped on rerun.

Per turn we log: parse success and the raw output on failure (format reliability); the
violation list from \texttt{check\_rules}, giving a violation taxonomy for free (rule
adherence, Table~\ref{tab:taxonomy}); and per-field $|\text{believed}-\text{true}|$ for
day, miles, food, bullets, and money (state tracking). In guided mode the true state is
shown every turn (can the model even copy correctly?); in blind mode it is shown only on
turn~1 and the model must reconstruct it from the event history. For the memory
experiment, three party facts stated once in the fixed game intro (Elsie's weak ankle,
Jonas's former trade, Ruben being the best shot) are quizzed at turns 5/10/15/20 at
temperature~0, scored by whether the gold answer appears in the reply. \texttt{analyze.py}
turns the log into the tables and figures below.

\begin{table}[t]
\caption{Hard-rule violation taxonomy audited by \texttt{engine.check\_rules}.}
\label{tab:taxonomy}
\centering
\scriptsize
\setlength{\tabcolsep}{3pt}
\begin{tabular}{lp{3.6cm}}
\toprule
Class & Meaning \\
\midrule
\texttt{miles\_exceeds\_max} & option moves $>25$ miles (or $<0$) \\
\texttt{overdraw:\emph{res}} & food/oxen/bullets/money pushed below 0 \\
\texttt{hunt\_*} & hunt with $<10$ bullets, without the $-10$ bullet cost, or yielding $>100$ lbs \\
\texttt{healing\_exceeds\_max} & one event heals $>40$ HP \\
\texttt{unknown\_party\_member} & effects name a nonexistent person \\
\texttt{recovery\_without\_health} & sick flag cleared at $\leq 30$ HP \\
\texttt{status\_on\_dead\_member} & flags set on a 0 HP member \\
\texttt{non\_numeric} / \texttt{non\_boolean} & wrong value types \\
\texttt{unknown\_sentiment} & sentiment outside the 6-word scale \\
\bottomrule
\end{tabular}
\end{table}

\section{Results}

\subsection{Format reliability and rule adherence}
Table~\ref{tab:format-rules} reports parse rate and violation density per strategy. (Turn
counts differ because games end early on death or arrival; runs lasted 12--20 turns.)
% Fill from results/tables.md ("Format reliability and rule adherence" table).
\begin{table}[t]
\caption{Format reliability and hard-rule violation rate per prompt strategy.}
\label{tab:format-rules}
\centering
\small
\begin{tabular}{lccc}
\toprule
Strategy & Turns & Parse rate & Viol.\ / 100 turns \\
\midrule
\texttt{minimal} & 66 & 0.985 & 184.6 \\
\texttt{fewshot\_schema} & 74 & 1.000 & 13.5 \\
\texttt{rules\_explicit} & 57 & 1.000 & 198.2 \\
\bottomrule
\end{tabular}
\end{table}

Format reliability is nearly solved even for \texttt{minimal}: 196 of 197 turns parsed,
the tolerant parser absorbing markdown fences and stray prose (the lone failure was an
\emph{unevaluated arithmetic expression}, \mbox{\texttt{"miles": 200 - (5 + 10 + \ldots)}},
foreshadowing the arithmetic weakness below). Rule adherence is not solved, and the
ordering is non-monotonic: \texttt{fewshot\_schema} commits $14\times$ fewer violations
than \texttt{minimal}, yet \texttt{rules\_explicit} is no better. The taxonomy explains it:
nearly all \texttt{minimal} and \texttt{rules\_explicit} violations are
\texttt{healing\_exceeds\_max}: the model echoing the party's \emph{absolute} health
values into the \emph{delta} field. One worked example with a negative delta
(\texttt{"Jonas": -10}) eliminates the class; stating the rule verbatim does not, because
the model does not know its numbers are read as deltas. A schema-\emph{interpretation}
error is fixed by an example, not a rule. As in the app, small models act on declarative
facts far more reliably than on conditionals they must evaluate themselves.

\subsection{State tracking and drift}
Guided error is $\sim$0 everywhere (Table~\ref{tab:state-tracking}): the model copies a
shown state faultlessly, so drift is not a transcription problem. Blind, belief diverges
immediately (Figure~\ref{fig:drift}): food is worst (MAE up to 129), miles reaches 43, yet
the day counter stays near-exact (0.8). The model tracks the calendar but loses any
quantity built from many small arithmetic updates, the failure the unevaluated-expression
parse error made explicit; the worked example halves food drift (56.4).
% Fill from results/tables.md ("State-tracking error" table).
\begin{table}[t]
\caption{Mean absolute error of the model's believed state vs.\ ground truth (parsed turns only).}
\label{tab:state-tracking}
\centering
\small
\setlength{\tabcolsep}{3.4pt}
\begin{tabular}{llccccc}
\toprule
Strategy & Mode & day & miles & food & bullets & money \\
\midrule
\texttt{minimal} & guided & 0.0 & 0.0 & 0.2 & 0.0 & 0.0 \\
\texttt{minimal} & blind & 0.8 & 42.9 & 126.9 & 6.2 & 35.5 \\
\texttt{fewshot} & guided & 0.0 & 0.5 & 0.6 & 0.0 & 0.0 \\
\texttt{fewshot} & blind & 0.8 & 22.8 & 56.4 & 3.2 & 36.2 \\
\texttt{rules} & guided & 0.0 & 0.0 & 0.0 & 0.0 & 0.0 \\
\texttt{rules} & blind & 0.8 & 28.4 & 129.0 & 4.8 & 37.2 \\
\bottomrule
\end{tabular}
\end{table}

\begin{figure}[t]
\centering
\IfFileExists{../results/drift_curves.png}{\includegraphics[width=\linewidth]{../results/drift_curves.png}}{\framebox[\linewidth]{\rule{0pt}{9em}\small run \texttt{run\_experiments.py} + \texttt{analyze.py}}}
\caption{State-tracking drift over the game in blind mode: mean $|\text{believed}-\text{true}|$ per turn for food (solid) and miles (dashed), per strategy.}
\label{fig:drift}
\end{figure}

One drift source was our own bug: food consumption (5 lbs per living member per day) was
originally charged silently by the harness, so blind-mode food drifted by the daily meal
cost regardless of recall; we were measuring a hidden rule, not context fidelity. Moving
the rule into the prompt removed the confound: a state-tracking experiment is only valid if
the state's dynamics are observable to the model.

\subsection{Memory of facts stated once}
Recall does not decay gradually; it falls off a cliff (Figure~\ref{fig:quiz}): 24/24
answers correct at turns 5 and 10, then 25\% at turn 15 and 6\% at turn 20, exactly when
the game intro (the facts' only source) leaves the 6{,}000-character history window. The
model was never \emph{remembering}, only \emph{reading}. This also confounds the
per-strategy scores in Table~\ref{tab:memory}: \texttt{rules\_explicit} looks best (0.850)
only because its games ended before the post-ejection quizzes. The lesson: durable facts
must be re-injected into every prompt, not left behind in the transcript.
% Fill from results/tables.md ("Memory recall" table).
\begin{table}[t]
\caption{Recall of party facts stated once at game start, quizzed at turns 5/10/15/20 (blind mode).}
\label{tab:memory}
\centering
\small
\begin{tabular}{lcc}
\toprule
Strategy & Quiz answers & Recall accuracy \\
\midrule
\texttt{minimal} & 32 & 0.500 \\
\texttt{fewshot\_schema} & 32 & 0.656 \\
\texttt{rules\_explicit} & 20 & 0.850 \\
\bottomrule
\end{tabular}
\end{table}

\begin{figure}[t]
\centering
\IfFileExists{../results/quiz_decay.png}{\includegraphics[width=\linewidth]{../results/quiz_decay.png}}{\framebox[\linewidth]{\rule{0pt}{9em}\small run \texttt{run\_experiments.py} + \texttt{analyze.py}}}
\caption{Fact recall vs.\ the turn at which the quiz was asked.}
\label{fig:quiz}
\end{figure}

\subsection{Qualitative findings from the app}
Three further errors surfaced in play, each with a fix. (1)~\emph{Narrative--arithmetic
incoherence}: one call doing both story and numbers produced consequences unrelated to its
prose, cured by the narrator/scorer split (Section~\ref{sec:split}). (2)~\emph{Option
monotony}: a schema example mentioning ``push on / wait / spend'' taught the model to offer
that trio every day, cured by demanding options differ \emph{in kind} plus scheduled
encounter days. (3)~\emph{Context leakage}: the 2-bullet scavenge cost surfaced in stories
as ``the party found 2 bullets'', cured by trimming redundant state from the prompt. A cap
on sentiment swings was reverted: it only fired when the model had good reason to swing,
manufacturing violations rather than catching them.

\section{Future Work}

Note that further optimizations could be done by either using VLLM or a smaller model where fit.
One area which has not been explored is running multiple models in parallel, where one model is used for the narration and another for the scoring.
One reason for this is that the behaviour of the inference engine plays a critical role in the performance when running multiple models in paralell.
The gameplay could be made more interesting by having an orchestrator prompt decide which event prompt to take and the event prompt being variable in itself, so that the model can be steered to a certain type of event.
I would end with noting how it is very phantomable given this experiment to image games containing LLMs as game masters or mere characters coming to the industry in the near feature.

\bibliography{references}
\bibliographystyle{icml2021}

\end{document}
